\documentclass[a4paper,12pt]{article}

\usepackage[utf8]{inputenc}
\usepackage[T1]{fontenc}
\usepackage[ngerman]{babel}
\usepackage{amsmath}
\usepackage{amssymb}
\usepackage{amsthm}
\usepackage{geometry}
\geometry{margin=2.5cm}

\theoremstyle{definition}
\newtheorem{definition}{Definition}[section]
\newtheorem{satz}[definition]{Satz}
\newtheorem{beispiel}[definition]{Beispiel}

\theoremstyle{remark}
\newtheorem{beweis}{Beweis}

\title{Mathematik -- Zusammenfassung\\
       \large Analysis I}
\author{Mathematikstudium}
\date{\today}

\begin{document}

\maketitle
\tableofcontents
\newpage

\section{Folgen und Konvergenz}

\begin{definition}[Konvergenz einer Folge]
Eine Folge $(a_n)_{n \in \mathbb{N}}$ heißt \textbf{konvergent} gegen den
Grenzwert $a \in \mathbb{R}$, wenn gilt:
\[
  \forall \varepsilon > 0 \;\exists N \in \mathbb{N} \;\forall n \geq N:
  \quad |a_n - a| < \varepsilon.
\]
Man schreibt dann $\lim_{n \to \infty} a_n = a$.
\end{definition}

\begin{satz}[Eindeutigkeit des Grenzwerts]
Der Grenzwert einer konvergenten Folge ist eindeutig bestimmt.
\end{satz}

\begin{beweis}
Seien $a$ und $b$ beide Grenzwerte von $(a_n)$. Für beliebiges
$\varepsilon > 0$ existieren $N_1, N_2 \in \mathbb{N}$ mit
$|a_n - a| < \varepsilon/2$ für $n \geq N_1$ und
$|a_n - b| < \varepsilon/2$ für $n \geq N_2$.
Für $n \geq \max(N_1, N_2)$ gilt dann:
\[
  |a - b| \leq |a - a_n| + |a_n - b| < \frac{\varepsilon}{2} +
  \frac{\varepsilon}{2} = \varepsilon.
\]
Da $\varepsilon > 0$ beliebig war, folgt $|a - b| = 0$, also $a = b$.
\end{beweis}

\begin{beispiel}
Die Folge $a_n = \dfrac{1}{n}$ konvergiert gegen $0$:
\[
  \lim_{n \to \infty} \frac{1}{n} = 0.
\]
\end{beispiel}

\section{Stetigkeit}

\begin{definition}[Stetigkeit]
Eine Funktion $f: D \to \mathbb{R}$ heißt \textbf{stetig} in $x_0 \in D$,
wenn gilt:
\[
  \forall \varepsilon > 0 \;\exists \delta > 0 \;\forall x \in D:
  \quad |x - x_0| < \delta \;\Rightarrow\; |f(x) - f(x_0)| < \varepsilon.
\]
\end{definition}

\begin{satz}[Zwischenwertsatz]
Sei $f: [a,b] \to \mathbb{R}$ stetig und $f(a) \cdot f(b) < 0$.
Dann existiert ein $c \in (a,b)$ mit $f(c) = 0$.
\end{satz}

\section{Differentialrechnung}

\begin{definition}[Ableitung]
Die Funktion $f: D \to \mathbb{R}$ heißt in $x_0 \in D$
\textbf{differenzierbar}, wenn der Grenzwert
\[
  f'(x_0) \;=\; \lim_{h \to 0} \frac{f(x_0 + h) - f(x_0)}{h}
\]
existiert. Man nennt $f'(x_0)$ die \textbf{Ableitung} von $f$ in $x_0$.
\end{definition}

\begin{satz}[Produktregel]
Sind $f$ und $g$ in $x_0$ differenzierbar, so ist auch $f \cdot g$
differenzierbar und es gilt:
\[
  (f \cdot g)'(x_0) = f'(x_0)\,g(x_0) + f(x_0)\,g'(x_0).
\]
\end{satz}

\end{document}
